% BOOK-FILE-ID: f-d581afb8e440
\subsection{非線形の周期性}\label{節：非線形の周期性}
一般項が求まらなくても,数列の性質を調べることはできる.
ここでは,数列が周期的になることを,一般項を求めずに示してみよう.

\paragraphbf{実験：2項の組が戻る}
$a_{n+2}=\frac{1+a_{n+1}}{a_n}$で初めの二項をともに$1$にすると,
$1,1,2,3,2,1,1,\ldots$となる.
第6,7項が初めの2項に戻った.次の項は直前の2項だけで決まるので,ここから先は同じ並びを繰り返すはずである.
では,初めの2項を変えても周期的になるのだろうか.正の$x,y$で確かめよう.

\paragraphbf{一般の初期値で確かめる}
\[
  \begin{aligned}
    a_1&=x>0,\qquad a_2=y>0,\\
    a_{n+2}&=\frac{1+a_{n+1}}{a_n}\quad(n\ge1)
  \end{aligned}
\]
とする.正の二項から次も正になるため,すべての項が定義される.

\paragraphbf{最初の2項への復帰}
\BookMathLayout{\[
    a_3=\frac{1+y}{x},\qquad
    a_4=\frac{1+x+y}{xy},\qquad
    a_5=\frac{1+x}{y}
  \]}{\[
    \begin{aligned}
      a_3&=\frac{1+y}{x},\\
      a_4&=\frac{1+x+y}{xy},\\
      a_5&=\frac{1+x}{y}
    \end{aligned}
  \]}
となる.$a_5$では
\[
  a_5=\frac{xy+1+x+y}{y(1+y)}=\frac{1+x}{y}
\]
と因数分解した.さらに
\[
  a_6=\frac{1+a_5}{a_4}=x,\qquad
  a_7=\frac{1+a_6}{a_5}=y
\]
と2項の組が戻り,その先も同じ漸化式で決まる.従って,
\[
  a_{n+5}=a_n\qquad(n\ge1)
\]
である.

\paragraphbf{周期と最小周期}
定\bookterm{sequence}{数列}になる$x=y=\frac{1+\sqrt5}{2}$では最小周期は1,ほかは5である.
定数解の条件は$x=y=t>0$, $t^2=t+1$から得る.
最小周期は5の約数である.そうでなければ,5を最小周期で割った正の余りも周期となり,最小性に反する.

\begin{bookpoint}
  \textbf{POINT}\par
  直前の2項から次が決まるなら,2項の組が元に戻ることを示せば周期性が言える.
\end{bookpoint}
ただし,初項からの周期,途中からの周期,符号だけの振動は区別すること.
\sectionref*{節：ガウス記号を含む式}の有限状態,\sectionref*{節：特性方程式の解が虚数の場合}の回転とは別の周期の仕組みである.
初項の例外は\bookproblemref{practice-connections}{5}{接続演習問5}で確かめよう.
